\documentclass[11pt,letterpaper]{article}
\usepackage[margin=.68in]{geometry}
\usepackage[T1]{fontenc}
\usepackage{lmodern,amsmath,amssymb,graphicx,booktabs,array,xcolor,fancyhdr,microtype,hyperref}
\definecolor{ink}{HTML}{173943}
\definecolor{muted}{HTML}{52646D}
\definecolor{light}{HTML}{EFF4F5}
\hypersetup{colorlinks=true,linkcolor=ink,urlcolor=ink,pdftitle={Mocap to matched camera images: a process guide},pdfauthor={Crazyflie camera research}}
\setlength{\parindent}{0pt}\setlength{\parskip}{6pt}
\setlength{\headheight}{14pt}
\pagestyle{fancy}\fancyhf{}\lhead{\small\color{ink}Crazyflie / 3D Gaussian splatting}\rhead{\small\color{muted}Process guide}\cfoot{\small\thepage}
\newcommand{\heading}[1]{\vspace{5pt}{\large\bfseries\color{ink}#1}\par}
\newcommand{\captiontext}[1]{{\small\color{muted}#1}\par}
\newcommand{\note}[1]{\vspace{3pt}\colorbox{light}{\parbox{\dimexpr\linewidth-2\fboxsep\relax}{\small #1}}\par}
\newcommand{\fig}[1]{\includegraphics[width=\linewidth]{docs/master_assets/#1}}
\begin{document}

{\LARGE\bfseries\color{ink}Mocap to matched camera images}\par
{\large A simplified explanation of the complete process}\par

The task is to make a reconstructed room appear from the same camera viewpoint
as a recorded Crazyflie image, then give the render a comparable camera response.
The room was captured at different times, so some objects and lighting differ.

\fig{overview.png}
\captiontext{One recorded frame and its simulation at the fitted viewpoint. Geometry and brightness are separate parts of the problem. Real pixels are preserved.}

\heading{The order matters}
\begin{center}\renewcommand{\arraystretch}{1.3}
\setlength{\tabcolsep}{3pt}
\begin{tabular}{@{}p{.285\linewidth}cp{.285\linewidth}cp{.285\linewidth}@{}}
\textbf{Synchronize time} & $\longrightarrow$ & \textbf{Find shared room points} & $\longrightarrow$ & \textbf{Correct every pose}\\
Pair images with mocap. && Associate real and rendered features. && Search camera rotation and position.
\end{tabular}
\end{center}
Then render the corrected poses, fit the simulated brightness/detail, and inspect
the images, overlays and separate error measures.

\heading{Start with the correct time and camera pose}
The camera clock and the motion-capture clock do not necessarily agree. Fit the
clock drift and offset, then test a remaining delay against the motion visible
between real images:
\[
t_q=t_{h,0}+a(t_d-t_{d,0})+b+\Delta.
\]
$t_d$ is camera-device time; $t_{h,0}$ is the first host timestamp; $a,b$ map the
centered device clock to host time; $\Delta$ is the remaining image-to-mocap delay.
$t_q$ is the time at which to query mocap. All terms use the same time unit.

Interpolate body position and orientation to $t_q$, then apply the calibrated
camera mount: $T_{MC}=T_{MB}T_{BC}$. Here $T_{AB}$ maps coordinates from frame $B$
into frame $A$: $M$ is mocap, $B$ is the tracked body and $C$ is the camera.
Missing tracking or large jumps are recorded as invalid poses.

\note{A recreated Motive rigid body can change its axes and pivot. Verify the camera-to-body transform whenever the tracked body or mount changes.}

\newpage
{\LARGE\bfseries\color{ink}Find the same physical features}\par

\heading{Use the same camera geometry}
Render with the Crazyflie's lens model so the simulated image has the same raw
pixel geometry as the recorded image. The real image is left unchanged.
For geometric calculations, invert the lens only at the measured feature points:
\[
p=\begin{bmatrix}(u-c_{0x})/f_x\\(v-c_{0y})/f_y\end{bmatrix}.
\]
$(u,v)$ is the point after lens inversion; $f_x,f_y$ are focal lengths in pixels;
$(c_{0x},c_{0y})$ is the camera principal point. $p$ is its normalized camera ray
coordinate. The active lens is an empirical calibration for the configured camera.

\heading{Bootstrap a usable starting alignment}
Render several known viewpoints of the room. Match real images to those views,
lift matched rendered pixels using scene depth, and estimate provisional camera
poses. These visual poses initialize the room transform and camera-axis correction.
Recheck timing after applying the camera correction.

\fig{feature_matches.png}
\captiontext{Matching letters/colors identify the same stored candidate association on both sides. Only a few candidates are drawn for clarity. These stored candidates illustrate the visual review step.}

\heading{Turn repeated observations into a fixed scene point}
SuperPoint detects image locations; LightGlue, with MINIMA weights for the
real/render pair, associates them. Match across rendered viewpoints, then
intersect the rays from their known camera poses:
\[
X_j=\arg\min_X\sum_k\left\|(I-d_kd_k^T)(X-p_k)\right\|^2.
\]
$X$ is a candidate point and $X_j$ the recovered room point; $p_k$ is the rendering camera centre, $d_k$ the
unit viewing ray and $I$ the identity matrix. The expression measures how far a
candidate point lies from all its viewing rays.

Keep points that agree across views, have useful viewing-angle separation and
reproject consistently in real images. Review the associations visually and
reject changed objects or repeated-pattern mistakes. Tracking propagates an
association through nearby frames only when consistency checks pass.

\note{Feature identity means the same physical location, not merely a similar-looking corner. Floor, gate and other scene-specific exclusions keep the room fit focused on stable background structure.}

\newpage
{\LARGE\bfseries\color{ink}Search the camera pose frame by frame}\par

\heading{First: one global initialization}
Use observations of admitted room points to establish the common scene registration. The
initial cameras receive a shared rotation, translation and scale correction:
\[
R_i'=R_cR_i,\qquad
p_i'=e^sR_c(p_i-c)+c+\delta,\qquad R_c=\exp([\omega]_\times).
\]
$R_i,p_i$ are the starting camera orientation and position; $\omega$ describes
the correction rotation; $\delta$ is its translation; $e^s$ is the distance
multiplier. $[\omega]_\times v=\omega\times v$ defines the rotation's skew matrix.
$c$ is the mean starting camera centre, used as a rotation/scale pivot.

\heading{Then: an independent correction for each observed frame}
Starting from the global fit, each frame has its own rotation $\omega_i$ and translation $\delta_i$:
\[
R_i'=\exp([\omega_i]_\times)R_i,\qquad p_i'=p_i+\delta_i.
\]
Project every known room point $X_j$ into the candidate camera pose and compare
its predicted position $\hat u_{ij}$ with the observed feature $u_{ij}$.
Both are in the same lens-corrected pixel coordinates.

\[
L(\theta)=\frac{1}{W}\sum_jw_j\,
\rho\!\left(\frac{\|\hat u_{ij}-u_{ij}\|}{\sigma_j}\right)
+\frac{1}{2W}\sum_k\left(\frac{\theta_k}{\tau_k}\right)^2.
\]
$\theta$ contains the candidate pose correction. $w_j$ is an observation's
weight, $W=\sum_jw_j$, and $\sigma_j$ is its pixel uncertainty. $\tau_k$ sets a
weak allowance for each correction component. The Huber penalty $\rho(e)$ is
$e^2/2$ for $e\leq1$, and $e-1/2$ otherwise: large feature errors have less influence.

\textbf{Search loop:} evaluate a coarse grid of candidate poses; retain the best
candidates; search smaller grids around them; finish with local refinement of
the same loss. Repeat this process using each frame's own observations.

\fig{pose_overlays.png}
\captiontext{Real image blended with a starting render and a per-frame render of the same frame. Doubled room edges help reveal a remaining pose mismatch. Use fixed room edges to assess the pose independently of brightness.}

\note{The global fit is only the starting point. Frames with no usable observations can receive a nearby/interpolated correction; longer unsupported gaps retain the global pose. The dataset labels these cases separately from independently fitted frames.}

\newpage
{\LARGE\bfseries\color{ink}Make the simulated appearance comparable}\par

Freeze the fitted poses before changing appearance. Convert the RGB render to
grayscale, then fit a shared response to the real recordings. This separates
camera/lighting differences from geometric misalignment.

\fig{appearance_dark.png}
\captiontext{An example from a darker recording. The middle panel changes geometry only; the right panel applies the frozen appearance model.}
\vspace{5pt}
\fig{appearance_bright.png}
\captiontext{An example looking toward bright light. Saturation and missing reconstructed detail remain visible; the response cannot recover information that was not captured.}

\[
\hat I(p)=\operatorname{clip}_{[0,1]}\!\left[
b+a(eh)^\beta\,\bar g(p)^\gamma
\exp\{-v_1q(p)-v_2q(p)^2+D(p)\}\right].
\]
\begin{tabular}{@{}p{.20\linewidth}p{.76\linewidth}@{}}
$\bar g(p)$ & Grayscale simulated intensity at pixel $p$, with a small selected blur.\\
$a,b,\gamma$ & Shared brightness gain, black offset and tone-curve exponent.\\
$e,h,\beta$ & Relative exposure times digital gain, a known-run correction, and the fitted exposure-response exponent.\\
$q(p),v_1,v_2$ & Squared image radius and the radial brightness coefficients.\\
$D(p)$ & Directional response from the scene-facing ray, plus a frame-level term from viewing direction and simulated bright content.
\end{tabular}

Fit the bounded coefficients using shared-feature pixels, uniformly distributed
tile averages, radial brightness and dark-image statistics. Select blur and
model form on development frames, then freeze the choice before checking
reserved frames. Tile averages balance where the loss is measured while one shared response
corrects the image.

\note{At inference, the correction requires the render, its camera orientation and recording settings. The empirical response approximates the combined appearance of scene lighting and the camera.}

\newpage
{\LARGE\bfseries\color{ink}Measure image and feature agreement}\par

\heading{Keep geometric error and image error separate}
For feature reprojection distances $e_j=\|\hat u_j-u_j\|$, report
\[
\mathrm{RMS}_{\rm feature}=\sqrt{\frac{1}{N}\sum_j e_j^2},
\qquad \mathrm{median}_{\rm feature}=\operatorname{median}_j e_j.
\]
$N$ is the number of compared observations. RMS emphasizes large misses;
the median describes a typical observation. Both measure feature agreement in image coordinates.

For image residuals $d_p=\hat I(p)-I(p)$ over the fixed comparison region
$\Omega$, use
\[
\mathrm{MAE}=\frac{1}{|\Omega|}\sum_{p\in\Omega}|d_p|,\qquad
\mathrm{RMSE}=\sqrt{\frac{1}{|\Omega|}\sum_{p\in\Omega}d_p^2}.
\]
$I$ is the real image and $\hat I$ is the simulated image, both on the same
intensity scale. RMSE is the square root of mean squared error. Signed mean
residual reveals whether the render is systematically brighter or darker.

\heading{Measure structure and edge agreement}
\textbf{SSIM} compares local brightness, contrast and image structure.\par
\textbf{Edge precision/recall} measure whether rendered and real edges occur
near each other. Combine them as $F_1=2PR/(P+R)$, where $P$ is precision and
$R$ is recall.\par
Use the separate components in a predeclared selection objective that also
penalizes radial brightness bias. This discourages a very dark or blurred
render from winning solely through lower pixel error. Report full-image and
stable-feature-region comparisons separately.

\fig{final_example.png}
\captiontext{The review product: real image, appearance-corrected simulation and their overlay. Look at fixed edges, black level, contrast and highlights independently.}

\heading{Read the outputs}
The pipeline saves synchronized pose validity, reviewed feature associations,
per-frame poses and fit modes, raw-geometry renders, corrected grayscale images,
and videos with an explicit frame/time mapping. These allow every step to be
checked against the original footage.

Moved objects, reconstruction gaps, glare, saturation and calibration error can
remain. Timing, lens, camera mount and feature identity need independent checks;
an appearance fit cannot repair them. Reserved images from the same recordings
test appearance conditional on the fitted poses, not transfer to a new room or
an independently calibrated camera.

\note{This guide explains the method with recorded examples. Reproduction commands, coefficient files and the full mathematical details are in the repository README and method guide.}

\end{document}
